\documentclass[12pt]{article}
\usepackage{amsmath, amsfonts, amssymb, amsthm, mathdots, mathrsfs, dsfont,systeme}
\usepackage[utf8]{inputenc}
\usepackage[spanish]{babel}
\usepackage{color}
\usepackage{fancyhdr}
\usepackage{enumerate}
\usepackage{hyperref}
\hypersetup{
    colorlinks=true,
    linkcolor=blue,
    filecolor=magenta,   
    citecolor=blue,
    urlcolor=cyan,
    pdftitle={Algebra lineal},
    pdfpagemode=FullScreen,
    }
\usepackage{amscd}
\usepackage{nomencl}
\usepackage{tikz}
\usepackage{tikz-cd}
\usepackage{amssymb}
\usepackage{blkarray}
\usepackage{multicol}
\usepackage[a4paper,margin=1cm,footskip=.5cm]{geometry}
\usepackage[all]{xy}
\usepackage{pst-node}
\usepackage{subfigure}
%\usepackage[table]{xcolor}
\usepackage{colortbl}
\newcommand{\cN}{\mathcal{N}}
\newcommand{\V}{\mathbb{V}}
\newcommand{\A}{\mathbb{A}}
\newcommand{\M}{\mathcal{M}}
\newcommand{\cH}{\mathcal{H}}
\newcommand{\F}{\mathcal{F}}
\newcommand{\G}{\mathcal{G}}
\newcommand{\cL}{\mathcal{L}}
\newcommand{\cO}{\mathcal{O}}
\newcommand{\Rn}{\mathbb{R}^n}
\DeclareMathOperator{\Hom}{Hom}
\DeclareMathOperator{\Ns}{NS}
\DeclareMathOperator{\im}{im}
\DeclareMathOperator{\Pic}{Pic}
\DeclareMathOperator{\Div}{Div}
\DeclareMathOperator{\Pdiv}{PDiv}
\DeclareMathOperator{\ddiv}{div}
\DeclareMathOperator{\Cl}{Cl}
\DeclareMathOperator{\CaDiv}{CaDiv}
\DeclareMathOperator{\Mor}{Mor}
\DeclareMathOperator{\Spec}{Spec}
\DeclareMathOperator{\id}{id}
\DeclareMathOperator{\Aut}{Aut}
\DeclareMathOperator{\End}{End}
\DeclareMathOperator{\CH}{CH}
\DeclareMathOperator{\IM}{Im}
\DeclareMathOperator{\Pf}{Pf}
\DeclareMathOperator{\Gl}{GL}
\DeclareMathOperator{\Sl}{SL}
\DeclareMathOperator{\St}{St}
\DeclareMathOperator{\tr}{Tr}
\DeclareMathOperator{\Gr}{Gr}
\DeclareMathOperator{\diag}{diag}
\DeclareMathOperator{\ord}{ord}
\DeclareMathOperator{\supp}{supp}
\DeclareMathOperator{\sign}{sign}
\DeclareMathOperator{\adj}{adj}
\DeclareMathOperator{\nul}{nul}
\DeclareMathOperator{\ran}{ran}
\newcommand{\calP}{\cal{P}}
\newcommand{\p}{\mathbb{P}}
\newcommand{\R}{\mathbb{R}}
\newcommand{\Z}{\mathbb{Z}}
\newcommand{\Q}{\mathbb{Q}}
\newcommand{\C}{\mathbb{C}}
\newcommand{\N}{\mathbb{N}}
\newcommand{\bbH}{\mathbb{H}}
\setlength{\parindent}{0pt}
\setlength{\textwidth}{6.5in}
\setlength{\textheight}{9.0in}
\setlength{\topmargin}{0in}
\setlength{\oddsidemargin}{0in}
\newtheoremstyle{exampstyle}
  {\topsep} % Space above
  {\topsep} % Space below
  {\itshape} % Body font
  {} % Indent amount
  {\bfseries} % Theorem head font
  {.} % Punctuation after theorem head
  {.5em} % Space after theorem head
  {} % Theorem head spec (can be left empty, meaning `normal')
\theoremstyle{exampstyle}
\newtheorem{thm}{Teorema}[section]
\newtheorem{prop}[thm]{Proposici\'on}
\newtheorem{lemma}[thm]{Lemma}
\newtheorem{cor}[thm]{Corolario}
\theoremstyle{definition}
\newtheorem{thme}{Expected Theorem}

\newtheorem{rec}[thm]{Recuerdo}
\newtheorem{conj}{Conjecture}
\newtheorem{defin}[thm]{Definici\'on}

\newtheoremstyle{break}%
{}{}%
{}{}%
{\bfseries}{}% % Note that final punctuation is omitted.
{\newline}{}
\theoremstyle{break}
\newtheorem{ejem}[thm]{Ejemplo}
\newtheorem{ejer}[thm]{Ejercicio}
\newtheorem{sol}[thm]{Soluci\'on}
\newtheorem{rmk}[thm]{Observaci\'on}
\numberwithin{equation}{section}
\newcommand{\gb}[1]{\operatorname{\textbf {#1}}}
\hypersetup{linkcolor=black}
\title{\'Algebra Lineal}
\author{Prof. Sebasti\'an Cea \\Prof. Camila Muñoz\\
}
\date{}
\pagestyle{fancy}

\lhead{\includegraphics[width = 4cm]{UFT.png}}
\setlength{\headsep}{0.6in}
\newcolumntype{a}{ >{\columncolor{yellow}} c }

\fancyhead[RO,LE]{
Prof: C. Muñoz.\\
01/2023}}



\begin{document}
%\maketitle
%\newpage

\begin{minipage}[r]{15cm}


\end{minipage}

\begin{center}
\Large{\textbf{\'Algebra Lineal\\Taller evaluado 2  }}
\end{center}

\section{Ejercicios}

\begin{enumerate}


\item Determine si el siguiente conjunto 
\[W_1= \left\{ \begin{pmatrix} x \\ y \end{pmatrix} \in \R{2} | x=3y \right\}\]
es un subespacio vectorial de $\R^{2}$.



\item Considere el conjunto de vectores $S = \left\{ \begin{pmatrix} 1 \\ 0 \\ 0 \end{pmatrix} , \begin{pmatrix} 2 \\ 1 \\ 0 \end{pmatrix} , \begin{pmatrix} 1 \\ 1 \\ 1 \end{pmatrix} \right\}$ en $\R^{3}$. Encuentre una base ortonormal de $\R^{3}$ a partir de $S$ usando el método de Gram-Schmidt.
    \item Encuentre la proyección de $\begin{pmatrix} 2 \\ 1 \\ 1 \end{pmatrix}$ sobre el plano $<\left\{ \begin{pmatrix} 1 \\ 0 \\ 0 \end{pmatrix} , \begin{pmatrix} 2 \\ 1 \\ 0 \end{pmatrix} \right\}  >$.


\item Determine si $\left\{ \begin{pmatrix} 1 \\ 2 \\ 1 \\ 1 \end{pmatrix}, \begin{pmatrix} 0 \\ 2 \\ 3 \\ 1 \end{pmatrix}, \begin{pmatrix} 0 \\ 0 \\ -3 \\ 0 \end{pmatrix}, \begin{pmatrix} 0 \\ 0 \\ 0 \\ 1 \end{pmatrix} \right\}$ es base de $\R^{4}$.

\item Dado $x \in \R^{n}$ un vector cualquiera, considere el conjunto
$$X^{\bot} = \{ y \in \R{n} | < y,x >=0 \}. $$
\begin{enumerate}
    \item Muestre que $X^{\bot}$ es subespacio vectorial de $\R^{n}$.
    \item En el caso en que $x= \Vec{0}$, ¿Cuál sería la dimensión de $X^{\bot}$?
    \item Determine la dimensión de $X^{\bot}$ cuando $x \neq \Vec{0}$.
\end{enumerate}
 


\end{enumerate}
\section*{Instrucciones}
\begin{itemize}
    \item Grupos de entre 4 y 5 personas.
    \item Se entrega al finalizar la clase.
\end{itemize}



\end{document}